\section{Описание}
Требуется для набора горизонтальных отрезков и набора точек определить, сколько отрезков лежит строго над каждой точкой.
Отрезок задан тройкой $(l, r, h)$ и содержит все точки $(x, h)$ при $l \le x \le r$. Точка $(x, y)$ лежит под отрезком,
если $l \le x \le r$ и $h > y$. Ограничения: $1 \le n, m \le 10^5$, координаты по модулю не превосходят $10^9$.
Запросы должны обрабатываться онлайн, то есть по одному, после построения структуры данных по всем отрезкам.

Задача сводится к двум условиям: отрезок должен накрывать абсциссу $x$ и его высота должна быть больше $y$.
Если обрабатывать запросы в произвольном порядке, нужно отвечать на вопрос \enquote{сколько отрезков,
накрывающих $x$, имеют высоту больше $y$}, причём для любого $x$ и в любом порядке. Для этого
используется \textbf{персистентное дерево отрезков}.

\subsection*{Идея решения}
Применим метод заметающей прямой по оси $x$. Каждый отрезок $[l, r]$ порождает два события:
\begin{itemize}
\item в точке $l$ отрезок \enquote{открывается}: в дерево добавляется $+1$ в позицию, соответствующую высоте $h$;
\item в точке $r + 1$ отрезок \enquote{закрывается}: в ту же позицию добавляется $-1$.
\end{itemize}
Закрытие в $r + 1$ (а не в $r$) превращает отрезок в полуоткрытый $[l, r + 1)$. Благодаря
этому отрезок учитывается и при $x = l$, и при $x = r$, а порядок событий с одинаковой координатой
$x$ не имеет значения: для запроса используются все события с координатой $\le x$.

Высоты предварительно сжимаются: после сортировки и удаления дубликатов каждой высоте соответствует номер
от $0$ до $k - 1$, где $k \le n$ --- число различных высот. Дерево отрезков строится над индексами
высот и хранит в каждой вершине количество \enquote{открытых} отрезков с высотами из соответствующего диапазона.

События сортируются по $x$. После применения каждого события создаётся новая версия дерева,
поэтому версия номер $v$ описывает состояние после первых $v$ событий; версия $0$ --- начальное дерево.
При изменении копируются только вершины на пути от корня до листа, остальные поддеревья используются
совместно со старой версией. Так версия создаётся за $O(\log n)$ времени и памяти, а все
предыдущие остаются доступными.

\subsection*{Обработка запроса}
Для запроса $(x, y)$:
\begin{enumerate}
\item бинарным поиском (\texttt{upper\_bound}) находится число событий $v$ с координатой $\le x$; нужная версия --- \texttt{roots[v]};
\item бинарным поиском по сжатым высотам находится номер $p$ первой высоты, строго большей $y$; если такой высоты нет, ответ равен $0$;
\item в версии \texttt{roots[v]} вычисляется сумма на суффиксе индексов $[p, k - 1]$: при спуске к листу $p$ каждый раз,
когда идём влево, прибавляется значение правого сына.
\end{enumerate}

\subsection*{Сложность}
Пусть $n$ --- число отрезков, $m$ --- число запросов.
\begin{itemize}
\item Сортировка событий и высот: $O(n \log n)$.
\item Построение всех версий: $2n$ событий по $O(\log n)$ вершин на каждое, то есть $O(n \log n)$ времени и памяти.
\item Один запрос: два бинарных поиска и один спуск по дереву, то есть $O(\log n)$. Итого на все запросы $O(m \log n)$.
\end{itemize}
Память: примерно $2n \cdot (\lceil \log_2 n \rceil + 1) \approx 3{,}6 \cdot 10^6$ вершин по $12$ байт, то есть около $43$ МБ.

\pagebreak

\section{Исходный код}
Программа состоит из следующих частей.

Вершина дерева отрезков хранится в структуре \texttt{Node}: индексы левого и правого сына и количество отрезков в поддереве.
Все вершины лежат в одном массиве \texttt{pool} (\texttt{std::vector<Node>}), а ссылки между вершинами представлены
индексами типа \texttt{int}, а не указателями. Это уменьшает размер вершины с $24$ до $12$ байт и снимает вопрос ручного
освобождения памяти: она освобождается автоматически при завершении программы.

Функция \texttt{createNode} строит начальное (нулевое) дерево. Функция \texttt{add} создаёт новую версию: копирует вершины
на пути к листу, а второго сына на каждом шаге берёт из старой версии без копирования. Результат рекурсивного вызова сначала
сохраняется в локальную переменную и только потом записывается в \texttt{pool}, потому что внутри вызова вектор может
перевыделить память, и ранее полученная ссылка на элемент стала бы недействительной. Функция \texttt{count\_height} считает
сумму на суффиксе индексов высот.

Структура \texttt{Event} хранит координату $x$ события, высоту отрезка и изменение $\pm 1$. В функции \texttt{main} читаются
отрезки, формируются и сортируются события, сжимаются высоты, строятся версии дерева и затем по одному обрабатываются запросы.

\begin{lstlisting}[language=C++]
#include <algorithm>
#include <iostream>
#include <iterator>
#include <vector>

struct Node {
    int left = -1;
    int right = -1;
    int count = 0;
};

std::vector<Node> pool;

int createNode(int l, int r) {
    int id = pool.size();
    pool.push_back(Node());
    if (l == r) {
        return id;
    }
    int mid = (l + r) / 2;
    int left = createNode(l, mid);
    int right = createNode(mid + 1, r);
    pool[id].left = left;
    pool[id].right = right;
    return id;
}

int add(int old_node_id, int l, int r, int pos, int delta) {
    int new_node_id = pool.size();
    pool.push_back(Node());
    pool[new_node_id].count = pool[old_node_id].count + delta;
    if (l == r) {
        return new_node_id;
    }
    int mid = (l + r) / 2;
    int left, right;
    if (pos <= mid) {
        left = add(pool[old_node_id].left, l, mid, pos, delta);
        right = pool[old_node_id].right;
    } else {
        left = pool[old_node_id].left;
        right = add(pool[old_node_id].right, mid + 1, r, pos, delta);
    }
    pool[new_node_id].left = left;
    pool[new_node_id].right = right;
    return new_node_id;
}

int count_height(int node_id, int l, int r, int pos) {
    int c = 0;
    int mid = (l + r) / 2;
    while (l < pos) {
        if (pos <= mid) {
            r = mid;
            c += pool[pool[node_id].right].count;
            node_id = pool[node_id].left;
        } else {
            l = mid + 1;
            node_id = pool[node_id].right;
        }
        mid = (l + r) / 2;
    }
    return c + pool[node_id].count;
}

struct Event {
    int x, h, delta;
};

int main(void) {
    int n, m;

    std::cin >> n >> m;

    std::vector<Event> events;
    std::vector<int> heights;
    events.reserve(2 * n);
    for (int i = 0; i < n; ++i) {
        int l, r, h;
        std::cin >> l >> r >> h;
        events.push_back({l, h, 1});
        events.push_back({r+1, h, -1});
        heights.push_back(h);
    }

    std::sort(
        events.begin(),
        events.end(),
        [](const auto& a, const auto& b) {return a.x < b.x;}
    );

    std::sort(heights.begin(), heights.end());
    heights.erase(
        std::unique(heights.begin(), heights.end()),
        heights.end()
    );

    std::vector<int> roots;
    roots.reserve(2 * n + 1);
    roots.push_back(createNode(0, n-1));
    for (auto& event : events) {
        auto it = std::lower_bound(heights.begin(), heights.end(), event.h);
        roots.push_back(add(
            roots.back(),
            0,
            n-1,
            std::distance(heights.begin(), it),
            event.delta
        ));
    }

    for (int i = 0; i < m; ++i) {
        int x, y;
        std::cin >> x >> y;

        auto it_event = std::upper_bound(
            events.begin(),
            events.end(),
            x,
            [](const auto& val, const auto& b){return val < b.x;}
        );
        auto v = std::distance(events.begin(), it_event);

        auto it_height = std::upper_bound(heights.begin(), heights.end(), y);
        if (it_height == heights.end()) {
            std::cout << 0 << '\n';
            continue;
        }
        auto pos = std::distance(heights.begin(), it_height);

        std::cout << count_height(roots[v], 0, n-1, pos) << '\n';
    }
}
\end{lstlisting}

В таблице приведено краткое описание основных элементов программы.
\begin{longtable}{|p{7.5cm}|p{7.5cm}|}
\hline
\rowcolor{lightgray}
\multicolumn{2}{|c|} {main1.cpp}\\
\hline
struct Node & Вершина дерева отрезков: индексы сыновей и счётчик\\
\hline
struct Event & Событие заметающей прямой: координата, высота, изменение $\pm 1$\\
\hline
int createNode(int l, int r) & Строит пустое дерево на отрезке индексов $[l, r]$, возвращает индекс корня\\
\hline
int add(int old\_node\_id, int l, int r, int pos, int delta) & Создаёт новую версию дерева с прибавлением \texttt{delta} в позиции \texttt{pos}, возвращает индекс нового корня\\
\hline
int count\_height(int node\_id, int l, int r, int pos) & Возвращает сумму значений на индексах $[pos, r]$ в данной версии\\
\hline
int main() & Чтение данных, построение версий, обработка запросов\\
\hline
\end{longtable}
\pagebreak

\section{Консоль}
Запуск на примере из условия задачи:
\begin{alltt}
$ g++ -pedantic -Wall -Wextra -std=c++17 -O2 -o sol main1.cpp
$ cat test1
5 3
0 5 2
1 3 1
1 8 5
-2 0 3
-1 6 -1
1 0
3 -3
3 5
$ ./sol < test1
3
4
0
\end{alltt}

Разберём первый запрос $(1, 0)$: точку $x = 1$ накрывают отрезки с высотами $2$, $1$, $5$ и $-1$ (отрезок $[-2, 0]$
высоты $3$ её не накрывает). Строго выше $y = 0$ лежат отрезки с высотами $2$, $1$ и $5$, поэтому ответ равен $3$.
Для запроса $(3, -3)$ точку накрывают отрезки высот $2$, $1$, $5$ и $-1$, и все они выше $-3$, ответ $4$.
Для запроса $(3, 5)$ нет отрезка выше высоты $5$, ответ $0$.

\pagebreak
